\documentclass[10pt,a4paper]{amsart}
\usepackage[margin=19mm]{geometry}
\usepackage{amsmath,amssymb,mathtools}
\usepackage[scr=rsfs,cal=euler]{mathalfa}
\usepackage{xfrac}
\usepackage[unicode,hidelinks,bookmarks=false]{hyperref}
\newcommand{\ie}{i.e.\ }
\newcommand{\cf}{cf.\ }
\newcommand{\resp}{respectively\ }
\newcommand{\Subsection}{\subsection}
\providecommand{\nobreakdash}{\nobreak\hbox{-}}
\setlength{\emergencystretch}{2em}
\multlinegap0pt
\usepackage{algorithm}\usepackage{algpseudocode}
\usepackage{float}
\newtheorem{lemma}{Lemma}
\newcommand{\ZZ}{\mathbb{Z}}
\title{A More Efficient Algorithm for Finding the Number of Permutations of $\ZZ/n\ZZ$ with Distinct Partial Sums}
\author{Baker, Quinn}
\date{Source: arXiv:2607.21892v1 (CC BY 4.0)}
\begin{document}
\maketitle

\noindent\emph{Source and licence.} A More Efficient Algorithm for Finding the Number of Permutations of $\ZZ/n\ZZ$ with Distinct Partial Sums. Version arXiv:2607.21892v1 (CC BY 4.0), distributed by arXiv under the Creative Commons Attribution 4.0 International licence, http://creativecommons.org/licenses/by/4.0/, as recorded in the arXiv licence metadata for this exact version and shown on the abstract page https://arxiv.org/abs/2607.21892v1. Full licence notice: \texttt{licenses/arxiv-2607.21892-CC-BY-4.0.txt}.\par\smallskip

\noindent\textit{Editorial scope.} Counted unit: the article's lemma lem:branch (source0.tex lines 175--179). The article's complete original proof of it is delivered with it (source0.tex lines 182--216, one \texttt{proof} environment of 35 lines). No mathematics is written, repaired or re-derived: the statement and the proof are the article's own lines, in the article's own notation, and no retained line is substituted or re-flowed.  Retained uncounted as the source-local context the proof needs: lines 142--143.  Nothing was omitted from any retained span, so the package carries no omission record.  No retained line contains a citation, so the wrapper carries no bibliography and no bibliography entry.  No retained line contains a figure environment, an image-inclusion command or drawing code, so no image is delivered and the package declares no asset dependency.  Editorial formatting only, in addition: an AMS \texttt{amsart} preamble that restates the article's own declarations and macros used by the retained lines, namely the environments \texttt{lemma} on separate counters, together with the standard packages those lines need.  The retained environments print in this document as Lemma 1 in order of appearance, whereas the article prints the same items with the numbering of its own full document.  The complete native source of the article as distributed in the arXiv e-print for this exact version is bundled unchanged as \texttt{sources/205899/source0.tex}.
\begin{lemma}\label{lem:rel_prime_constructive} Let $\sigma$ be any permutation of $\ZZ / n\ZZ$ and let $r$ be any element of $\ZZ / n\ZZ$ such that $\gcd(r,n)=1$. Then $(\sigma_0,\sigma_1,\ldots,\sigma_{n-1})$ is a sequencing if and only if $(r\sigma_0, r\sigma_1,\ldots,r\sigma_{n-1})$ is.
\end{lemma}

\begin{lemma}\label{lem:branch}
Let $C$ be the set of all sequencings of $\ZZ / n\ZZ$. Partition $C$ into branches of sequencings $C_i$ with $1\leq i\leq n-1$, where
\[C_i:= \{(0,i,g_2,\ldots,g_{n-1})\in C \}\]
and let $D$ be the set of all proper divisors of $n$. Then \[\#C = \sum_{d\in D}\#C_d\cdot \varphi(n/d).\]
\end{lemma}

\begin{proof}

For each $d\in D$ define
%\[
%M_d = \bigcup_{j=1}^{n-1}\bigl\{ g \mid g \textnormal{ is the smallest element of } (\ZZ/n\ZZ)^\times \textnormal{ with }dg\equiv j\pmod{n}\bigr\}.
%\]

\[
M_d = \{1 \leq i \leq n/d \mid (i,n)=1\}.
\]

It is a well-established fact that $\varphi(n/d) = \#\{1 \leq a \leq n \mid (a,n)=d\}$, and thus we find that
\begin{align*}
\varphi(n/d) &= \lvert\{1 \leq a \leq n \mid (a,n)=d\}\rvert \\
&=  \lvert  \{ 1 \leq a/d \leq n/d \mid (a/d,n/d)=1\}\rvert \\
&= \lvert \{1 \leq i \leq n/d \mid (i,n)=1\}\rvert \\
&=\# M_d.
\end{align*}


Let \[W:=\bigcup_{d\in D} \bigcup_{g\in M_d} (C_d \cdot g).\]

Then $\#W=\sum_{d\in D}\#C_d\cdot \varphi(n/d)$. Thus we will prove this lemma by proving $W=C$.

We first establish that $W\subseteq C$. This follows immediately from Lemma \ref{lem:rel_prime_constructive} which demonstrates that multiplying a sequencing by an element relatively prime to $n$ still gives a sequencing. 

We next show that $C\subseteq W$. Pick any $t\in C$. Then $t$ is a sequencing $c\in C$ and write $c=\{0,g_1,g_2,\ldots,g_{n-1}\}.$ The first nonzero element of $c$, $g_1$ can be factored into $g_1\equiv g\cdot d$ (mod $n$) for some $d\in D$. 

As $g$ is relatively prime to $n$, its inverse $g^{-1}\in\ZZ/n\ZZ$ exists, is relatively prime to $n$ and $g_1g^{-1}\equiv d\ $ (mod $n$). Now consider
$\{0,g_1g^{-1},g_2g^{-1},\ldots,g_{n-1}g^{-1}$.
This  is a sequencing and it starts with $d$ as the first nonzero element so $\{0,g_1g^{-1},g_2g^{-1},\ldots,g_{n-1}g^{-1}\}\in C_d$. Therefore, \[c=\{0,g_1g^{-1},g_2g^{-1},\ldots,g_{n-1}g^{-1}\}\cdot g.\] As this is the form of all elements of $W$ we have that $c\in W$ and thus $C\subseteq W$.

Therefore we have $C=W$.

\end{proof}

\end{document}
